\documentclass[11pt]{ctexart}
\usepackage[a4paper,top=2.2cm,bottom=2.2cm,left=2.2cm,right=2.2cm]{geometry}
\usepackage{booktabs}
\usepackage{longtable}
\usepackage{array}
\usepackage{enumitem}
\usepackage{xcolor}
\usepackage{amssymb}
\usepackage[hidelinks]{hyperref}
\setlist{nosep, leftmargin=1.6em}
\emergencystretch=1.5em
\setlength{\tabcolsep}{4pt}
\renewcommand{\arraystretch}{0.92}

\hypersetup{
  pdftitle={SmartFarmAgent 智慧农业智能体系统——项目开题方案（导师提交版）},
  pdfsubject={项目开题方案（导师提交版，v1.5）},
  pdfauthor={}
}

\begin{document}

\begin{center}
  {\LARGE\bfseries SmartFarmAgent 智慧农业智能体系统}\\[4pt]
  {\large 项目开题方案（导师提交版）}
\end{center}

\vspace{0.5em}

\section{项目背景与痛点}
\begin{itemize}
  \item 传统大棚管理依赖人工巡检与经验调控，存在缺水/涝害、病害发现滞后、水肥浪费等问题。
  \item 市面「智慧农业」多为阈值规则控制，无真正的智能决策；大模型应用停留在「知识问答」层，没有闭环到设备执行。
  \item 2025—2026 年竞赛集中在「智能体 / IoT + 大模型」方向；农业开源大模型与 RAG 框架已成熟可用，硬件成本低，整条链路可自研。
\end{itemize}

\section{系统总体架构}
\begin{longtable}{p{1.2cm}p{2.2cm}p{8.8cm}p{3.0cm}}
  \toprule
  层级 & 名称 & 组成与职责 & 关键技术 \\
  \midrule
  L6 & 展示层 & 微信小程序 / H5 大屏 / LabVIEW 上位机 & 前端、上位机 \\
  L5 & 智能决策层 & 大模型智能体（LLM+RAG）→ JSON 设备指令（核心创新）；视觉服务（YOLO 病害/生长分析） & LLM、RAG、视觉 \\
  L4 & 云端平台层 & MQTT Broker（EMQX/OneNET）→ 时序存储 → Node-RED 可视化 & MQTT、时序存储 \\
  L3 & 边缘层 & ESP32-CAM 拍照/轻量推理；断网规则兜底（本地优先） & 边缘计算、轻量推理 \\
  L2 & 通信层 & USART 文本帧协议 + ESP8266 MQTT & 串口、MQTT \\
  L1 & 感知执行层 & STM32F103C8T6 + 传感器 + 继电器/泵/灯/喷药 & GPIO、ADC、I2C、SPI \\
  \bottomrule
\end{longtable}
安全原则：本地阈值规则永远优先于云端指令（断网自持、误操作兜底），云端智能体指令走白名单。

\section{功能模块}
\begin{longtable}{p{2.2cm}p{5.4cm}p{0.9cm}p{4.4cm}}
  \toprule
  功能 & 实现方案 & 难度 & 可行性依据 \\
  \midrule
  自动浇水/施肥 & 土壤湿度+光照综合判断；继电器驱动水泵/蠕动泵 & 中 & 传感器与继电器已有；开源整机项目已验证 \\
  环境自动调控 & DHT11/BH1750 → 风扇/加热/补光灯；PWM 调光（PID 可选） & 中 & 传感器已有；PWM/PID 为已学内容 \\
  病虫害检测+给药 & 摄像头定拍 → YOLOv8（先云后边）→ 置信度达标 → 喷药泵 & 高 & PlantDoc/PlantVillage 公开数据集 + 多个 ESP32-CAM/YOLO 先例 \\
  生长状况分析 & 定距定时拍照 → 叶片数/株高/颜色统计 → 大模型生成周报 & 高 & 视觉测量工具包开源可用 \\
  智能体决策 & 传感器+检测+历史 → 大模型（司农/DeepSeek）→ 设备指令+解释 & 高 & Dify HTTP 节点可调外部 API；农业大模型已开源 \\
  数据中台/展示 & MQTT 上行 → 时序存储 → 小程序/H5；W25Q64 本地日志 & 中 & EMQX/Node-RED 成熟；农业岛平台可二开 \\
  \bottomrule
\end{longtable}

\section{数据链路}
\begin{longtable}{p{1.5cm}p{11.0cm}p{2.6cm}}
  \toprule
  环节 & 内容 & 输出 \\
  \midrule
  感知 & BH1750/DHT11/土壤 → STM32 采集打包（文本帧）；摄像头（ESP32-CAM）→ 图片上传 & 数据帧 / 图片 \\
  通信 & USART 帧 + ESP8266 MQTT（本地同时接 LabVIEW 调试 + W25Q64 日志） & 上行数据 \\
  云端 & MQTT Broker（EMQX/OneNET）→ 时序存储 + Node-RED 可视化 & 时序数据 \\
  AI & 视觉服务（YOLO 病害检测/生长分析）→ 结果入事件队列；智能体（LLM+RAG）综合传感器+视觉+历史 → 决策 JSON & 决策 JSON \\
  执行 & MQTT 下行指令 → 网关转 HTTP/串口 → STM32 解析 → 继电器/水泵/风扇/补光灯/喷药泵 & 设备动作 \\
  闭环 & 执行状态回传 → 数据入库 → 小程序/大屏展示 + 智能体下次决策依据 & 反馈数据 \\
  \bottomrule
\end{longtable}

\section{技术可行性}
\begin{enumerate}
  \item 硬件与嵌入式：全部器件为 STM32 生态常用模块，与现有学习内容一一对应，参考整机项目已跑通相同硬件组合。
  \item 视觉检测：PlantDoc（427 星）/ PlantVillage（5 万+ 图、38 类）公开数据集可用，YOLOv8 框架成熟（Ultralytics 60k 星），已有温室病害检测与 ESP32-CAM 实时检测先例。
  \item 智能体链路：司农（国内首个农业开源大语言模型，8B/32B）与 AgriAgent（125 星）均已开源；Dify 官方确认 HTTP 节点可调外部 API，「智能体 → HTTP → MQTT 网关 → 设备」路径可行。
  \item 云平台：EMQX（16.5k 星）、Node-RED（23.5k 星）、农业岛（347 星）均已核验；OneNET 免费可用。
  \item 已知约束：ESP32-CAM 算力有限，初期视觉推理走云端；司农 32B 本地部署显存要求高，初期用 8B 或 DeepSeek API。
\end{enumerate}

\section{里程碑计划}
\begin{longtable}{p{2.0cm}p{1.5cm}p{7.0cm}p{4.1cm}}
  \toprule
  阶段 & 时间 & 内容 & 验收点 \\
  \midrule
  M0 开题 & 第 1—2 周 & 链路定稿、软件环境搭建（Python/Dify/YOLO） & 开题通过 \\
  M1 感知-控制闭环 & 第 3—6 周 & 传感器采集 → OLED → 继电器 → 串口帧协议 → LabVIEW 上位机 & 本地闭环运行，数据帧正确，远程手动可控 \\
  M2 上云+日志 & 第 7—10 周 & ESP8266 MQTT → EMQX/OneNET → Node-RED 可视化；W25Q64 日志 & 云平台实时曲线、断网日志完整可补传 \\
  M3 视觉检测 & 第 11—18 周 & 公开集+自采数据训练 YOLO（先 5 类常见病）→ 云端推理 → 触发喷药 & 检测准确率 $\geq 85\%$，喷药动作联动 \\
  M4 智能体决策 & 第 19—23 周 & Dify 搭智能体+RAG 知识库 → 综合决策 → 指令闭环 → 小程序 & 自然语言指令自动执行，误操作有兜底 \\
  M5 作品化 & 第 24 周 & 外壳/文档/演示视频/说明书/答辩材料 & 达到参赛交付标准 \\
  \bottomrule
\end{longtable}

\section{创新点与预期成果}
\begin{enumerate}
  \item 智能体闭环决策：大模型直接输出设备动作指令，不只停留在问答层。
  \item 边缘-云协同：断网时本地规则独立运行，联网后云端智能体接管；本地规则永远优先，误操作有兜底。
  \item 低成本：硬件总成本低，农户可负担，便于推广。
  \item 全周期数字档案：从种子到收获持续记录传感器与视觉数据，形成可追溯的生长档案。
\end{enumerate}

\section{风险与应对}
\begin{longtable}{p{5.2cm}p{9.5cm}}
  \toprule
  风险 & 应对 \\
  \midrule
  视觉识别数据不足/效果差 & 公开数据集预训练 + 自采微调；范围先缩到 3—5 类常见病害 \\
  ESP32-CAM 边缘算力不足 & 初期视觉推理放云端，边缘只拍照上传；后期再试轻量模型 \\
  智能体误操作 & 本地规则优先、指令白名单、关键动作人工确认模式 \\
  时间不足 & 按里程碑裁剪：合并施肥/给药、生长分析降级为拍照+周报 \\
  \bottomrule
\end{longtable}

\section{主要开源参考}
\href{https://github.com/langgenius/dify}{Dify}（智能体/工作流平台，151k 星）、\href{https://github.com/infiniflow/ragflow}{RAGFlow}（RAG 引擎，87k 星）、\href{https://github.com/ultralytics/ultralytics}{Ultralytics}（YOLO 框架，60k 星）、\href{https://github.com/pratikkayal/PlantDoc-Dataset}{PlantDoc}（病害数据集，427 星）、PlantVillage（\href{https://www.kaggle.com/datasets/abdallahalidev/plantvillage-dataset}{Kaggle}，5 万+ 图、38 类）、司农（南京农业大学农业大语言模型，8B/32B）、\href{https://github.com/zhiweihu1103/AgriAgent}{AgriAgent}（农业多模态大模型，125 星）、\href{https://github.com/emqx/emqx}{EMQX}（16.5k 星）/ \href{https://github.com/node-red/node-red}{Node-RED}（23.5k 星）（MQTT Broker 与可视化工具）。

\end{document}
